\documentclass[a4paper,11pt]{report}
\usepackage[utf8]{inputenc}
\usepackage[vietnamese]{babel}
\usepackage{amsmath,amssymb}
\usepackage{mathptmx}
\usepackage[margin=2.2cm]{geometry}
\usepackage{booktabs}
\usepackage{tabularx}
\usepackage{ragged2e}
\usepackage{enumitem}
\usepackage{microtype}
\usepackage{listings}
\usepackage{xcolor}
\usepackage[hidelinks]{hyperref}

% Cấu hình cột bảng chống giãn khoảng cách trắng
\newcolumntype{Y}{>{\RaggedRight\arraybackslash}X}
\newcolumntype{L}[1]{>{\RaggedRight\arraybackslash}p{#1}}

% Cấu hình listings cho code snippet gọn gàng
\lstset{
  basicstyle=\ttfamily\small,
  breaklines=true,
  frame=single,
  backgroundcolor=\color{gray!5},
  rulecolor=\color{gray!40},
  keywordstyle=\bfseries\color{blue!70!black},
  commentstyle=\itshape\color{green!50!black},
  stringstyle=\color{red!60!black},
  showstringspaces=false
}

\begin{document}

% TRANG BÌA GIÁO TRÌNH
\begin{titlepage}
    \centering
    \vspace*{1.5cm}
    
    {\Huge \textbf{GIÁO TRÌNH KỸ NGHỆ VIBECODING CHUYÊN SÂU}\par}
    \vspace{0.8em}
    {\LARGE \textit{Agentic Software Engineering, Open-Source Arsenal \& Hacker Scripting}\par}
    
    \vspace{1.5cm}
    \rule{0.8\textwidth}{1.5pt}\par
    \vspace{1.5cm}
    
    {\Large \textbf{Từ Tư Duy Xếp Hình Lego, Giải Phẫu Vũ Khí Nguồn Mở Đến Scripting Chuyên Sâu Linux/Windows Và Vòng Lặp Tự Phục Hồi}\par}
    
    \vfill
    
    {\large \textbf{Trung Tâm Nghiên Cứu \& Thực Hành Vibecoding Antigravity}\par}
    \vspace{0.5em}
    {\normalsize Dựa trên 28 Công trình Nghiên cứu Thực nghiệm 2026 Đã Thẩm Định Bằng Playwright\par}
    \vspace{1.5cm}
    
    {\large Tháng 10 Năm 2026}\par
\end{titlepage}

% LỜI NÓI ĐẦU
\chapter*{Lời Nói Đầu}
\addcontentsline{toc}{chapter}{Lời Nói Đầu}

Sự bùng nổ của trí tuệ nhân tạo tạo sinh (\textit{Generative AI}) đã đưa ngành công nghiệp phần mềm bước vào một bước ngoặt lịch sử: thuật ngữ \textbf{Vibecoding} ra đời, mô tả một trạng thái mà ở đó con người chỉ cần đưa ra ý tưởng bằng ngôn ngữ tự nhiên và để các mô hình ngôn ngữ lớn (LLM) gánh vác toàn bộ việc sinh mã. 

Tuy nhiên, sau sự hào nhoáng ban đầu, thực tế khắc nghiệt nhanh chóng lộ diện: đa số các ứng dụng do người mới phát triển đều thất bại hoặc không thể đưa vào sử dụng thực tế. Ứng dụng thường mắc phải căn bệnh ``nhồi nhét tính năng vô tội vạ'', giao diện mang vẻ ngoài ``rất nhựa'' sáo rỗng (\textit{AI Slop}), mã nguồn chứa đầy các lỗ hổng bảo mật nghiêm trọng (hơn 65\% code AI dính lỗi CWE theo nghiên cứu 2026), và sau vài lần sửa lỗi, toàn bộ hệ thống bị tê liệt do nợ nhận thức tích lũy (\textit{Cognitive Debt}).

Bộ giáo trình này được biên soạn nhằm giải quyết tận gốc rễ những vấn đề trên. Đây không phải là một cuốn sổ tay lý thuyết chung chung. Đây là một \textbf{Giáo trình Kỹ nghệ Thực chiến Toàn diện} được thiết kế để đào tạo bạn từ một người dùng thụ động trở thành một **Tổng Kiến Trúc Sư kiêm Hacker Ghép Nối Lego**:
\begin{enumerate}[label={\arabic*.}]
    \item \textbf{Chương 1}: Khởi nguyên, bản chất nhận thức học và triết lý Lego Hacker: Đứng trên vai những người khổng lồ.
    \item \textbf{Chương 2}: Phân tích nghiệp vụ bằng Jobs-to-be-Done (JTBD), MoSCoW và Quy tắc Cắt 50\% tính năng râu ria.
    \item \textbf{Chương 3}: Giải phẫu Tech Stack thực chiến: làm chủ Playwright, công cụ thay thế tối ưu (Crawl4AI, HTTPX), sức mạnh toàn diện của Python và vòng lặp tự sửa lỗi (*Self-Healing loop*).
    \item \textbf{Chương 4}: An ninh mã nguồn: phòng thủ Top 5 lỗ hổng bảo mật CWE nguy hiểm nhất do AI sinh ra.
    \item \textbf{Chương 5}: Kiến trúc Logic \& Schema-First: máy trạng thái 4 bước và quy tắc mô-đun siêu nhỏ dưới 200 dòng.
    \item \textbf{Chương 6}: Thiết kế UX/UI Chống Nhựa (Anti-AI-Slop): tỷ lệ màu 60--30--10 và nhịp điệu khoảng trắng 8px.
    \item \textbf{Chương 7}: Chuyên sâu Scripting \& OS Internals: Nghệ thuật đường ống Unix (Linux) đối đầu Object Pipeline (Windows PowerShell) và vòng lặp tự chủ khép kín.
\end{enumerate}

Chúng tôi tin rằng, khi nắm vững tư duy kiến trúc và làm chủ các công cụ kiểm soát, bạn sẽ khai phóng được toàn bộ tiềm năng của trí tuệ nhân tạo mà vẫn đảm bảo tính an toàn, thẩm mỹ cao cấp và sự trường tồn của sản phẩm.

\vspace{1.5em}
\noindent\textbf{Ban Biên Soạn Giáo Trình Vibecoding Antigravity}

% MỤC LỤC
\tableofcontents

% CHƯƠNG 1
\chapter{Khởi Nguyên, Bản Chất Vibecoding \& Triết Lý Lego Hacker}

\section{Lịch sử hình thành: Từ Autocomplete đến Kỷ nguyên Tự chủ Agentic}
Thuật ngữ \textit{``Vibecoding''} được khởi xướng bởi Andrej Karpathy vào đầu năm 2025, nhằm mô tả một phương thức lập trình mà ở đó con người hoàn toàn không trực tiếp gõ từng dòng mã nguồn, mà chỉ giao tiếp bằng ngôn ngữ tự nhiên, thả lỏng vào ``dòng chảy cảm xúc'' (\textit{vibes}) và để mô hình AI đảm nhận toàn bộ việc hiện thực hóa.

Để hiểu đúng bản chất kỹ thuật, chúng ta cần đặt Vibecoding vào tiến trình lịch sử 4 giai đoạn của công nghệ hỗ trợ lập trình:
\begin{enumerate}[label={\arabic*.}]
    \item \textbf{Giai đoạn 1: Dự đoán cục bộ (Autocomplete \& Copilot 2021--2023)}: Các mô hình ngôn ngữ hoàn thành câu lệnh dựa trên ngữ cảnh ngắn cục bộ (vài chục dòng mã). Con người vẫn là người viết chính, mô hình đóng vai trò gợi ý phím Tab.
    \item \textbf{Giai đoạn 2: Hỗ trợ đối thoại (Chat-Assisted Pair Programming 2023--2024)}: Các giao diện trò chuyện cho phép dán lỗi hoặc yêu cầu sinh hàm độc lập. Nhược điểm chí mạng là con người phải sao chép qua lại bằng tay và thiếu ngữ cảnh toàn cục của dự án.
    \item \textbf{Giai đoạn 3: Tích hợp môi trường phát triển (IDE-Integrated Agents 2024--2025)}: Sự xuất hiện của Cursor, Windsurf với khả năng lập chỉ mục cơ sở mã nguồn, đọc tệp cục bộ và chỉnh sửa nhiều tệp (\textit{multi-file editing}).
    \item \textbf{Giai đoạn 4: Tác tử tự chủ vòng lặp kép (Autonomous Agentic Engineering 2025--2026)}: Các hệ thống như Antigravity (Google DeepMind) và SWE-agent thực thi các chu kỳ khép kín: Tự phân tích $\rightarrow$ Tự gọi công cụ (Terminal, LSP, Browser/Playwright) $\rightarrow$ Tự chạy kiểm thử $\rightarrow$ Tự phát hiện ngoại lệ $\rightarrow$ Tự sửa lỗi cho đến khi kiểm thử vượt qua hoàn toàn (\textit{loop-until-green}).
\end{enumerate}

\section{Triết lý Lego Hacker: Đứng trên vai những người khổng lồ}
Người lập trình nghiệp dư thường cố gắng viết lại mọi thứ từ đầu. Ngược lại, một **Lego Hacker** hiểu rằng thế giới mã nguồn mở đã tạo ra những vũ khí CLI hoàn hảo nhất:
\begin{itemize}[leftmargin=1.5em]
    \item Muốn tải video và phụ đề từ 1.000 trang web? Dùng \textbf{yt-dlp}.
    \item Muốn nén, ghép, lọc luồng video/âm thanh? Dùng \textbf{FFmpeg}.
    \item Muốn chiếu và điều khiển điện thoại Android? Dùng \textbf{scrcpy} và \textbf{ADB}.
    \item Muốn chuyển đổi tài liệu Word/Excel sang PDF? Dùng \textbf{LibreOffice Headless} và \textbf{Pandoc}.
    \item Muốn nhận dạng giọng nói offline? Dùng \textbf{Whisper}.
\end{itemize}

Công thức tối thượng:
\begin{center}
\fbox{\parbox{0.88\textwidth}{
\textbf{Công thức sản phẩm Vibecode siêu tốc:}\\
$\text{Sản phẩm Hoàn Thiện} = \text{Giao diện / Kịch bản mỏng} + \text{Vũ khí Nguồn Mở có sẵn} + \text{AI làm chất keo kết nối}$
}}
\end{center}

\section{Bóc trần ảo giác: Nợ nhận thức và bẫy tái triển khai}
Đa số người mới tiếp cận Vibecoding đều rơi vào một bẫy tâm lý nguy hiểm: \textbf{Ảo giác làm chủ (Illusion of Competence)}. Nghiên cứu thực nghiệm của \textit{Xu và cộng sự (arXiv:2609.09022, 2026)} đã chỉ ra:
\begin{itemize}[leftmargin=1.5em]
    \item \textbf{Nợ nhận thức tích lũy (Cognitive Debt Accumulation)}: Phát sinh khi tốc độ sinh mã của AI vượt xa tốc độ đọc hiểu và kiểm chứng của con người, tạo ra cơ sở mã mì ăn liền sụp đổ dây chuyền.
    \item \textbf{Bẫy tái triển khai (Reimplementation Trap - arXiv:2609.29410)}: LLM tự viết lại các hàm tiện ích đã tồn tại trong dự án thay vì import các thư viện nguồn mở có sẵn.
\end{itemize}

\section{Quy Trình Tạo Mẫu Siêu Tốc: AI Studio Apps, OpenCode và 9Router}
Trong thực tế phát triển sản phẩm năm 2026, quy trình vibecoding hiệu quả nhất thường đi theo mô hình kết hợp giữa đám mây và tác tử cục bộ:
\begin{enumerate}[label={\arabic*.}]
    \item \textbf{Google AI Studio Apps}: Động cơ tạo mẫu MVP tức thì (\texttt{https://aistudio.google.com/apps}) sử dụng Gemini 2.0/2.5 cho phép sinh giao diện sống tương tác được trong 30 giây.
    \item \textbf{OpenCode (\texttt{anomalyco/opencode})}: Sau khi tải mã nguồn ZIP về, sử dụng tác tử dòng lệnh TUI viết bằng Go (chế độ \texttt{plan} và \texttt{build}) để refactor lại code đơn khối, bóc tách API về backend và loại bỏ cạm bẫy lộ API Key ở máy khách.
    \item \textbf{9Router (\texttt{9router.com})}: Đóng vai trò là Gateway thông minh đa nhà cung cấp (\texttt{http://localhost:20128/v1}), hỗ trợ tự động rẽ nhánh 3 tầng (\textit{3-Tier Auto-Fallback}) khi gặp lỗi giới hạn hạn mức (HTTP 429), kết hợp bộ nén Rust Token Killer (RTK) giúp tiết kiệm 20\% đến 65\% chi phí token.
\end{enumerate}

\section{Cây quyết định: Lựa chọn hướng đi và mô hình Vibecoding}

\begin{table}[htbp]
\centering
\small
\caption{Ma trận quyết định lựa chọn lộ trình Vibecoding thực chiến}
\vspace{0.5em}
\begin{tabularx}{\textwidth}{L{3.2cm} L{3.8cm} L{4.5cm} Y}
\toprule
\textbf{Lộ trình thực chiến} & \textbf{Công cụ phối hợp nòng cốt} & \textbf{Đặc điểm kiến trúc} & \textbf{Đối tượng \& Thời gian tối ưu} \\
\midrule
\textbf{Lộ trình A: Thần tốc MVP (24h Ship)} & AI Studio Apps + OpenCode TUI + 9Router & React Vite / Next.js, Server Actions, SQLite nhúng, Tailwind CSS & Indie Hacker, kiểm chứng ý tưởng thị trường trong 1--2 ngày \\
\textbf{Lộ trình B: Kỹ sư Tự trị (Autonomous Swarm)} & Antigravity CLI / OpenCode Headless + Multi-Agent (Spec, Coder, Auditor, QA) & Schema-First, Zod contracts, TDD loop-until-green, Micro-modules & Ứng dụng SaaS phức tạp, Full-stack web, SDK, CLI tools \\
\textbf{Lộ trình C: Hacker Cục bộ (Local-First \& Offline)} & Local Ollama (DeepSeek R1/V3) + PocketBase / DuckDB + Bash/Python & 100\% Offline, không truyền dữ liệu ra ngoài, chi phí API bằng 0 & Xử lý dữ liệu nhạy cảm nội bộ, doanh nghiệp đề cao bảo mật \\
\bottomrule
\end{tabularx}
\end{table}

\noindent\textbf{Kịch bản thực hành đồng hành}: Tệp \texttt{token\_budget\_calculator.py} (Chương 1) giúp người học tính toán độ suy giảm chú ý của ngữ cảnh theo số lượng token và ước lượng chi phí API 2026.

% CHƯƠNG 2
\chapter{Phân Tích Nghiệp Vụ \& MVP Scoping: Khắc Phục Hội Chứng Nhồi Nhét Tính Năng}

\section{Bản chất bài toán: Ảo tưởng về sự phong phú}
Khi chi phí sinh mã nguồn giảm về gần bằng không nhờ trí tuệ nhân tạo, lập trình viên rất dễ mắc hội chứng nhồi nhét tính năng: yêu cầu AI làm đồng thời đăng nhập, dark mode, bảng xếp hạng, xuất file và diễn đàn. Kết quả là làm phân mảnh trải nghiệm người dùng, bùng nổ các trường hợp biên và làm cạn kiệt bộ nhớ chú ý của AI.

\section{Khung lý thuyết Jobs-to-be-Done (JTBD) trong Vibecoding}
Theo Giáo sư Clayton Christensen: \textit{``Khách hàng không mua sản phẩm, họ thuê sản phẩm để hoàn thành một bước tiến trong cuộc sống.''}

Công thức định vị JTBD chuẩn mực:
\begin{center}
\fbox{\parbox{0.85\textwidth}{
\textit{Khi gặp hoàn cảnh [Bối cảnh cụ thể], tôi muốn [Hành động trực tiếp], để tôi có thể đạt được [Kết quả chuyển đổi cốt lõi].}
}}
\end{center}

Ví dụ: Thay vì làm một ứng dụng luyện thi TOEIC rườm rà, JTBD cốt lõi là: \textit{``Khi gặp câu hỏi khó có nhiều từ mới, học sinh muốn bóc tách ngữ pháp đơn giản trước để làm chủ cấu trúc, sau đó mới tích hợp từ vựng học thuật, để không bị hoảng loạn khi làm bài.''} Sản phẩm chỉ cần một Tính năng Sát thủ (Killer Feature) duy nhất giải quyết bài toán này.

\section{Thang đo MoSCoW và Quy tắc Cắt 50\% Tính Năng}
Toàn bộ yêu cầu ban đầu phải được phân loại qua MoSCoW:
\begin{itemize}[leftmargin=1.5em]
    \item \textbf{Must Have}: Tính năng sống còn quyết định giá trị của sản phẩm.
    \item \textbf{Should Have}: Tính năng quan trọng nhưng có thể giải quyết thủ công tạm thời.
    \item \textbf{Could Have / Won't Have}: Tính năng râu ria, xao nhãng; kiên quyết loại bỏ khỏi phiên bản đầu tiên.
\end{itemize}
\textbf{Quy tắc Cắt 50\%}: Gạch bỏ ngay lập tức $50\%$ số tính năng phụ. Khóa chặt phạm vi vào: \textbf{1 Persona + 1 Core Loop + 1 North Star Metric}.

\section{Kỹ thuật Phỏng vấn Đối kháng Socratic (/grill-me) và spec.md}
Trước khi viết mã, hãy yêu cầu AI đóng vai một Giám đốc Sản phẩm khó tính để tra khảo ngược lại các giả định và trường hợp biên của bạn. Sau đó, đóng băng toàn bộ kết quả vào tệp đặc tả kỹ thuật máy đọc được (\texttt{spec.md}) với danh mục \textit{Non-Goals} minh bạch.

\section{Cây quyết định: Lựa chọn hướng đi sản phẩm sau phân tích nghiệp vụ}

\begin{table}[htbp]
\centering
\small
\caption{Cây quyết định phân loại hướng đi sản phẩm sau bước phân tích nghiệp vụ}
\vspace{0.5em}
\begin{tabularx}{\textwidth}{L{3.2cm} L{3.8cm} L{4.5cm} Y}
\toprule
\textbf{Hướng đi sản phẩm} & \textbf{Chiến lược phạm vi (Scope)} & \textbf{Kiến trúc kỹ thuật tối ưu} & \textbf{Chỉ số thành công (Metric)} \\
\midrule
\textbf{Hướng 1: Micro-SaaS Đơn Tính Năng} & 1 tính năng sát thủ duy nhất, không bắt đăng ký rườm rà, giải quyết tức thì & Next.js / Vite + Server Actions + Stripe / LemonSqueezy & Tỷ lệ chuyển đổi người dùng dùng thử thành công $> 35\%$ \\
\textbf{Hướng 2: Tiện Ích Cục Bộ (Local-First)} & Hoạt động 100\% offline, không phụ thuộc đám mây, bảo mật dữ liệu tuyệt đối & Python / Tauri + SQLite FTS5 + PocketBase nhúng & Tốc độ khởi động $< 1$s, dung lượng bộ cài $< 15$MB \\
\textbf{Hướng 3: Tự Động Hóa \& Bot Dòng Lệnh} & Không cần UI, tập trung xử lý luồng công việc lặp đi lặp lại & Python CLI (argparse/click) + Linux Bash / Windows PowerShell & Tiết kiệm thời gian thao tác thủ công từ hàng giờ về vài giây \\
\bottomrule
\end{tabularx}
\end{table}

\noindent\textbf{Kịch bản thực hành đồng hành}: Tệp \texttt{socratic\_spec\_builder.py} (Chương 2) hỗ trợ người học tự động hóa sinh tệp \texttt{spec.md} hoàn chỉnh qua 5 câu hỏi cốt lõi của JTBD.

% CHƯƠNG 3
\chapter{Giải Phẫu Tech Stack Thực Chiến: Playwright, Python \& Vòng Lặp Tự Sửa Lỗi}

\section{Playwright Thực Chiến và Bảng Đối Chiếu Công Cụ Tối Ưu}
Trong quy trình Vibecoding tác tử tự chủ, Playwright đóng vai trò là ``đôi mắt và bàn tay'' giúp mô hình AI tương tác trực tiếp với ứng dụng.

\subsection{Bảy siêu năng lực thực tế của Playwright}
\begin{enumerate}[label={\arabic*.}]
    \item \textbf{Headless Browser Automation}: Khởi chạy đa nhân Chromium, Firefox, WebKit ngầm không cần giao diện đồ họa.
    \item \textbf{Lắng nghe \& Bắt chặn gói tin mạng (Network Interception)}: Bắt trộm dữ liệu máy khách gửi lên (\texttt{page.on('request')}) và giả lập phản hồi API (\texttt{page.route()}) để kiểm thử mã lỗi 500 hoặc mạng chập chờn.
    \item \textbf{Cơ chế tự động chờ (Auto-waiting)}: Tự động chờ phần tử hiển thị và sẵn sàng nhận thao tác mà không cần dùng lệnh chờ thô sơ.
    \item \textbf{Bắt lỗi Console \& Ngoại lệ chưa xử lý}: Bắt toàn bộ \texttt{console.error} và \texttt{pageerror} (như React Hydration mismatch, Null pointer crash).
    \item \textbf{Chụp ảnh màn hình kiểm thử thị giác (Visual Diff)}: Chụp ảnh full-page làm minh chứng đối chiếu giao diện cho AI Agent.
    \item \textbf{Lưu trữ phiên đăng nhập (Storage State)}: Ghi nhận cookies và localStorage vào tệp \texttt{auth.json} để tái sử dụng.
    \item \textbf{Trích xuất DOM động từ Single Page Applications (SPA)}: Chờ ứng dụng React/Next.js hoàn tất render dữ liệu máy khách rồi mới trích xuất.
\end{enumerate}

\subsection{Khi nào Playwright quá tải? Các công cụ thay thế nhẹ hơn và dễ hơn}
Mỗi phiên Chromium ngốn từ 250MB đến 600MB RAM. Việc lạm dụng Playwright cho mọi tác vụ đọc web là một sai lầm nghiêm trọng gây tràn RAM hệ thống (\textit{OOM}).

\begin{table}[htbp]
\centering
\small
\caption{So sánh Playwright với các công cụ thu thập và tự động hóa thay thế}
\vspace{0.5em}
\begin{tabularx}{\textwidth}{L{2.5cm} L{2cm} L{2.2cm} L{3.5cm} Y}
\toprule
\textbf{Công cụ} & \textbf{Tiêu thụ RAM} & \textbf{Tốc độ} & \textbf{Đặc tính cốt lõi} & \textbf{Khi nào nên sử dụng?} \\
\midrule
\textbf{Playwright} & 300--600 MB & 2--5 giây & Mô phỏng đầy đủ Chromium, click DOM, bắt console log & Bắt buộc khi test WebApp tương tác, bắt lỗi console, chụp ảnh giao diện \\
\textbf{Crawl4AI} & 80--150 MB & 0.5--1.5 giây & Chuyên cho LLM, tự chuyển HTML thành Markdown sạch & Rất hay khi cào bài báo, blog, tài liệu để nhồi vào prompt của AI \\
\textbf{Browser-Use / Stagehand} & 400--700 MB & Chậm (gọi AI) & Tác tử tự nhìn màn hình và tự click bằng ngôn ngữ tự nhiên & Dùng khi web có cấu trúc quá phức tạp, selector CSS thay đổi liên tục \\
\textbf{HTTPX + BeautifulSoup} & \textbf{10--20 MB} & \textbf{$< 0.2$ giây} & Client HTTP thuần túy, không tốn tài nguyên trình duyệt & Rất tối ưu khi cào arXiv, Wikipedia, API tĩnh; tiết kiệm 95\% RAM \\
\bottomrule
\end{tabularx}
\end{table}

\section{Hai Quy Trình Thực Chiến Bằng Mã Nguồn Mẫu}

\subsection{Quy trình 1: Thu thập và Thẩm định Bài báo Nghiên cứu Tự động}
Thay vì mở trình duyệt nặng nề, kỹ sư sử dụng HTTPX gọi thẳng arXiv API để lấy cấu trúc dữ liệu XML chỉ trong 0.8 giây; chỉ khi cần chụp ảnh minh chứng thị giác mới khởi chạy Playwright ngầm:

\begin{lstlisting}[language=python, caption={Kich ban thu thap bai bao sieu nhe voi Playwright va HTTPX}]
import httpx, xml.etree.ElementTree as ET
from playwright.sync_api import sync_playwright

def fetch_arxiv_fast(paper_id: str):
    url = f"https://export.arxiv.org/api/query?id_list={paper_id}"
    resp = httpx.get(url, follow_redirects=True, timeout=10.0)
    root = ET.fromstring(resp.text)
    entry = root.find('{http://www.w3.org/2005/Atom}entry')
    title = entry.find('{http://www.w3.org/2005/Atom}title').text.strip()
    abstract = entry.find('{http://www.w3.org/2005/Atom}summary').text.strip()
    return {"id": paper_id, "title": title, "abstract": abstract}
\end{lstlisting}

\subsection{Quy trình 2: Tự Động Test WebApp \& Vòng Lặp Tự Sửa Lỗi (Self-Healing)}
Kịch bản lắng nghe toàn bộ 4 kênh lỗi: console.error, uncaught pageerror, network 4xx/5xx và requestfailed. Khi phát hiện lỗi, script ghi dữ liệu vào tệp máy đọc \texttt{agent\_feedback.json} kèm số dòng cụ thể và thoát với Exit Code 1. AI Agent đọc tệp này để tự khoanh vùng và sửa code cho đến khi Exit Code là 0:

\begin{lstlisting}[language=python, caption={Co che xuat phan hoi loi co cau truc cho AI Agent}]
# scripts/self_healing_test_runner.py (Core Implementation)
errors = []
page.on("console", lambda m: errors.append({"type": "console", "text": m.text, "loc": m.location}) if m.type == "error" else None)
page.on("pageerror", lambda e: errors.append({"type": "exception", "msg": str(e)}))
page.on("response", lambda r: errors.append({"type": "http", "status": r.status, "url": r.url}) if r.status >= 400 else None)

page.goto("http://localhost:3000", wait_until="networkidle")
if errors:
    page.screenshot(path="failure_screenshot.png")
    with open("agent_feedback.json", "w") as f:
        json.dump({"status": "FAILED", "errors": errors}, f, indent=2)
    sys.exit(1) # Exit Code 1: Agent continues fixing code
sys.exit(0)     # Exit Code 0: All green, loop completed!
\end{lstlisting}

\section{Sức Mạnh Toàn Diện Của Python Trong Kỷ Nguyên Vibecoding}
Python không chỉ là ngôn ngữ khoa học dữ liệu, mà trong Vibecoding nó chính là \textbf{chất keo tự động hóa tối thượng} với 6 trụ cột:
\begin{enumerate}[label={\arabic*.}]
    \item \textbf{Tự động hóa kịch bản \& File system}: Quản lý tiến trình, biên dịch tài liệu LaTeX tự động, dọn dẹp các tiến trình Chrome ma bị treo ngầm.
    \item \textbf{Cào dữ liệu \& Xử lý tri thức}: Tích hợp nhịp nhàng giữa HTTPX (nhẹ, nhanh) và Playwright (vượt JS) để nạp tài liệu cho AI.
    \item \textbf{Backend hiệu năng cao với FastAPI \& Pydantic}: Xây dựng API bất đồng bộ siêu tốc, tự động sinh tài liệu Swagger/OpenAPI và kiểm duyệt dữ liệu nghiêm ngặt.
    \item \textbf{Xử lý PDF \& Tài liệu chuyên sâu}: Bóc tách từng trang PDF đề thi bằng \texttt{pypdf}, \texttt{pdfplumber}, xuất công thức Toán KaTeX/LaTeX chuẩn mực.
    \item \textbf{Xử lý dữ liệu cấu trúc}: Dùng Polars và Pandas biến đổi hàng nghìn dòng dữ liệu JSON/CSV chỉ trong một phần nghìn giây.
    \item \textbf{Điều phối AI cục bộ \& Embeddings}: Kết nối các mô hình cục bộ qua LiteLLM, llama.cpp và tính toán vector ngữ nghĩa.
\end{enumerate}

\section{Đại Thư Viện Vũ Khí Nguồn Mở: Jina Reader, PocketBase, DuckDB và Cloudflared}
Để xây dựng sản phẩm thần tốc, kỹ sư Vibecoding cần làm chủ các vũ khí nguồn mở chuyên sâu:
\begin{itemize}[leftmargin=1.5em]
    \item \textbf{Jina Reader (\texttt{https://r.jina.ai/[URL]})}: Bóc tách bất kỳ trang web nào thành Markdown sạch sẽ cho LLM mà không cần khởi chạy trình duyệt, tiết kiệm 100\% RAM so với Playwright.
    \item \textbf{Nodriver}: Thư viện tự động hóa Chrome không dùng webdriver flag, giao tiếp trực tiếp qua Chrome DevTools Protocol (CDP) giúp vượt qua hoàn toàn các bộ lọc bot của Cloudflare Turnstile và Akamai.
    \item \textbf{PocketBase}: Tệp binary duy nhất viết bằng Go (chỉ 15MB) tích hợp SQLite, Authentication, File Storage S3, Realtime SSE và Admin UI, thay thế hoàn toàn nhu cầu dựng cụm Supabase hay Firebase cho các bản MVP.
    \item \textbf{DuckDB \& SQLite FTS5}: Cỗ máy SQL phân tích dữ liệu lớn trong bộ nhớ cực nhanh, kết hợp cùng FTS5 mang lại khả năng tìm kiếm toàn văn và tô đậm từ khóa chỉ với vài dòng code Python thuần.
    \item \textbf{Cloudflared Tunnel}: Tiện ích dòng lệnh chính thức của Cloudflare cho phép mở cổng \texttt{localhost:3000} ra Internet công khai dưới dạng tên miền HTTPS bảo mật miễn phí chỉ bằng 1 dòng lệnh (\texttt{cloudflared tunnel --url http://localhost:3000}) mà không cần mở port router hay sở hữu IP tĩnh.
    \item \textbf{Typst}: Công cụ soạn thảo văn bản khoa học thế hệ mới viết bằng Rust, biên dịch tài liệu sang PDF chỉ mất vài chục mili-giây, cú pháp trực quan thay thế sự cồng kềnh kéo dài của LaTeX cổ điển.
\end{itemize}

\section{Cây quyết định: Lựa chọn Tech Stack Thu Thập Dữ Liệu \& Lưu Trữ}

\begin{table}[htbp]
\centering
\small
\caption{Cây quyết định lựa chọn công cụ thu thập dữ liệu \& lưu trữ}
\vspace{0.5em}
\begin{tabularx}{\textwidth}{L{2.8cm} L{3.5cm} L{4.5cm} Y}
\toprule
\textbf{Nhu cầu bài toán} & \textbf{Công cụ khuyến nghị} & \textbf{Cơ chế hoạt động tối ưu} & \textbf{Hạn chế kỹ thuật} \\
\midrule
\textbf{Lấy dữ liệu API / Trang tĩnh} & \textbf{HTTPX} (Async) & Gửi HTTP GET/POST trực tiếp, chỉ tốn 15MB RAM, tốc độ $< 0.8$s & Bất lực trước SPA render bằng JS (React/Vue/Angular) \\
\textbf{Bóc tách trang web thành Markdown} & \textbf{Jina Reader} (\texttt{r.jina.ai}) & Thêm tiền tố \texttt{r.jina.ai/URL}, trả về Markdown sạch 0 dòng code & Phụ thuộc dịch vụ bên ngoài, giới hạn rate limit \\
\textbf{Vượt Bot Detection / Cloudflare} & \textbf{Nodriver} / \textbf{Crawl4AI} & Điều khiển Chrome DevTools Protocol (CDP) không gắn cờ webdriver & Cấu hình phức tạp hơn Playwright truyền thống \\
\textbf{Test Webapp \& Chụp màn hình} & \textbf{Playwright} & Bắt toàn diện 4 kênh lỗi (Console, PageError, Network 4xx/5xx) & Tốn 300--500MB RAM cho mỗi tiến trình trình duyệt \\
\midrule
\textbf{Lưu trữ Cục bộ / Đơn người dùng} & \textbf{SQLite} (WAL mode) & 1 tệp tin duy nhất, tốc độ đọc triệu dòng/s, tích hợp FTS5 & Nghẽn ghi khi có nhiều người dùng ghi đồng thời \\
\textbf{Phân tích Dữ liệu / OLAP} & \textbf{DuckDB} & Tối ưu truy vấn cột (\textit{columnar}), đọc trực tiếp Parquet/CSV & Không phù hợp cho ứng dụng web giao dịch (OLTP) \\
\textbf{Backend Tức thì / Auth + DB} & \textbf{PocketBase} & 1 file nhúng nhị phân duy nhất, có sẵn Auth, Storage và Realtime & Không mở rộng quy mô đa máy chủ ngang (horizontal scale) \\
\textbf{Production SaaS / Đa người dùng} & \textbf{PostgreSQL} (Supabase) & Đầy đủ RLS, ACID, PgBouncer, quan hệ khóa ngoại phức tạp & Cần cấu hình kết nối cẩn trọng để tránh cạn Connection Pool \\
\bottomrule
\end{tabularx}
\end{table}

\noindent\textbf{Kịch bản thực hành đồng hành}: Thư mục Chương 3 cung cấp \texttt{paper\_harvester.py} (thu thập arXiv siêu tốc) và \texttt{self\_healing\_test\_runner.py} (quét 4 kênh lỗi webapp).

% CHƯƠNG 4
\chapter{An Ninh Mã Nguồn Trong Kỷ Nguyên Vibecoding: Phòng Thủ Lỗ Hổng CWE \& OWASP}

\section{Thực trạng an ninh: Nghiên cứu thực nghiệm 2026}
Theo nghiên cứu \textit{Last Secure Code Benchmark (arXiv:2610.14766, 2026)}, hơn **65\% các đoạn mã chức năng do các mô hình LLM hàng đầu sinh ra chứa ít nhất một lỗ hổng bảo mật CWE nghiêm trọng**. Lý do là AI được tối ưu hóa để làm cho mã chạy được theo prompt, chứ không tự động thêm các tầng phòng thủ nếu không được yêu cầu.

\section{Top 5 lỗ hổng bảo mật chí mạng}
\begin{enumerate}[label={\arabic*.}]
    \item \textbf{CWE-89 (SQL Injection)}: AI nối chuỗi trực tiếp vào câu lệnh SQL thô (\texttt{queryRawUnsafe} hoặc f-strings). Khắc phục: Luôn dùng Parameterized Queries.
    \item \textbf{CWE-639 / CWE-862 (IDOR \& BOLA)}: Lấy tài nguyên chỉ dựa vào ID từ URL mà không kiểm tra xem \texttt{user\_id} có thuộc về người đăng nhập hay không. Khắc phục: Ràng buộc quyền sở hữu ngay trong mệnh đề WHERE.
    \item \textbf{CWE-284 (Bỏ quên RLS trong Supabase)}: AI tạo bảng nhưng không bật \texttt{ENABLE ROW LEVEL SECURITY}, làm lộ toàn bộ bảng dữ liệu qua \texttt{anon\_key}. Khắc phục: Luôn kích hoạt RLS và thiết lập chính sách Policy tường minh.
    \item \textbf{CWE-798 \& CWE-200 (Rò rỉ khóa bí mật)}: Gắn tiền tố \texttt{NEXT\_PUBLIC\_} vào API Secret Key, khiến khóa bị nhúng vào bundle JavaScript của trình duyệt. Khắc phục: Chỉ đọc khóa bí mật ở môi trường máy chủ.
    \item \textbf{CWE-918 (SSRF)}: Nhận URL của người dùng và gọi trực tiếp từ backend, làm lộ địa chỉ IP nội bộ và metadata đám mây. Khắc phục: Sử dụng danh sách tên miền tin cậy và chặn dải IP riêng tư.
\end{enumerate}

\section{Cây quyết định: Thiết lập thế trận phòng thủ an ninh mã nguồn}

\begin{table}[htbp]
\centering
\small
\caption{Cây quyết định thiết lập thế trận phòng thủ an ninh mã nguồn}
\vspace{0.5em}
\begin{tabularx}{\textwidth}{L{3.2cm} L{3.8cm} L{4.5cm} Y}
\toprule
\textbf{Thế trận phòng thủ} & \textbf{Đặc trưng dữ liệu} & \textbf{Hàng rào an ninh bắt buộc} & \textbf{Chi phí triển khai} \\
\midrule
\textbf{Thế trận 1: Công cụ Công khai (Public Read-Only)} & Dữ liệu mở, không có tài khoản người dùng cá nhân (ví dụ: tra cứu đề thi) & Rate Limiting (Cloudflare / Upstash), ẩn API Key ở Backend, chặn SSRF IP nội bộ & Thấp: thiết lập trong vòng 15 phút \\
\textbf{Thế trận 2: Tiện ích Cục bộ (Local-First Offline)} & Dữ liệu người dùng lưu trực tiếp trên máy (SQLite, IndexedDB) & Mã hóa tệp DB nếu nhạy cảm, không mở cổng mạng công khai, kiểm tra quyền file & Rất thấp: không có rủi ro tấn công qua mạng Internet \\
\textbf{Thế trận 3: Nền tảng SaaS Đa Khách Hàng (Multi-Tenant)} & Dữ liệu tài chính, người dùng khác nhau chia sẻ chung database & RLS Supabase bắt buộc, kiểm tra quyền sở hữu IDOR ở mọi endpoint, Cookie HttpOnly & Cao: đòi hỏi kiểm thử tự động Pentest trước khi deploy \\
\bottomrule
\end{tabularx}
\end{table}

\noindent\textbf{Kịch bản thực hành đồng hành}: Tệp \texttt{cwe\_security\_scanner.py} (Chương 4) tự động phát hiện 5 lỗ hổng CWE điển hình của AI và trả về Exit Code 1 phục vụ tự động sửa lỗi.

% CHƯƠNG 5
\chapter{Kiến Trúc Logic \& Schema-First: Chống Trôi Ngữ Cảnh và Máy Trạng Thái Tường Minh}

\section{Triết lý Schema-First và Hợp đồng dữ liệu}
Để ngăn chặn hiện tượng trôi ngữ cảnh (\textit{Context Drift}) khi AI tự ý đổi tên trường giữa các tệp, toàn bộ mô hình dữ liệu phải được khóa cứng trước bằng TypeScript Interfaces và Zod/Pydantic schemas. Mọi luồng dữ liệu đi vào hệ thống đều phải đi qua cổng kiểm duyệt (\textit{Boundary Gatekeepers}) bằng phương thức \texttt{safeParse}.

\section{Máy trạng thái tường minh 4 bước (Explicit 4-State Machine)}
Tuyệt đối xóa bỏ các cờ boolean rời rạc dễ gây xung đột (\texttt{isLoading}, \texttt{isError}). Thay vào đó, chuẩn hóa mọi component theo 4 trạng thái hữu hạn bất biến:
\begin{center}
\textbf{IDLE} $\longrightarrow$ \textbf{LOADING} $\longrightarrow$ \textbf{SUCCESS} $\longrightarrow$ \textbf{ERROR}
\end{center}
Phân định rạch ròi trạng thái \textit{Empty State} (thành công nhưng dữ liệu rỗng) và \textit{Error State} (lỗi ngoại lệ hoặc mất mạng).

\section{Định luật Mô-đun siêu nhỏ dưới 200 dòng mã}
Khi kích thước tệp vượt quá 300 dòng mã, tỷ lệ AI sinh lỗi hồi quy và làm mất code cũ tăng vọt. Hãy phân rã cơ sở mã thành các tệp chuyên biệt dưới 200 dòng: \texttt{[feature].types.ts}, \texttt{[feature].service.ts}, \texttt{[feature].ui.tsx}, và \texttt{[feature].test.ts}.

\section{Cơ chế Ngắt Mạch (Circuit Breaker)}
Cài đặt bộ ngắt mạch 3 trạng thái (\textit{Closed, Open, Half-Open}) khi gọi API bên ngoài để bảo vệ hệ thống không bị treo khi bên thứ ba gặp sự cố hoặc cạn hạn mức yêu cầu.

\section{Cây quyết định: Lựa chọn mô hình kiến trúc logic phần mềm}

\begin{table}[htbp]
\centering
\small
\caption{Cây quyết định lựa chọn mô hình kiến trúc logic phần mềm}
\vspace{0.5em}
\begin{tabularx}{\textwidth}{L{3.2cm} L{3.8cm} L{4.5cm} Y}
\toprule
\textbf{Mô hình kiến trúc} & \textbf{Nguyên lý tổ chức} & \textbf{Công nghệ nòng cốt} & \textbf{Quy mô \& Bối cảnh tối ưu} \\
\midrule
\textbf{Mô hình 1: Micro-Modules Monolith} & Đơn khối phân rã hạt nhân, mỗi tệp $< 200$ dòng, tách biệt schema / service / UI & Next.js / Vite SPA + Server Actions + Zod + SQLite/Supabase & MVP, ứng dụng đơn mục tiêu, quy mô dưới 50.000 người dùng \\
\textbf{Mô hình 2: Serverless Edge Microservices} & Luồng bất đồng bộ, truyền phát trực tiếp (SSE Streaming), ngắt mạch Circuit Breaker & Cloudflare Workers / Vercel Edge + Upstash Redis Queue & Chatbot AI, xử lý hàng đợi, ứng dụng có lưu lượng truy cập đột biến \\
\textbf{Mô hình 3: Local-First Event-Driven} & Xử lý cục bộ 100\%, đồng bộ ngầm khi có mạng, bảo mật tuyệt đối & Tauri / Electron + SQLite FTS5 / PGlite + CRDTs sync & Ứng dụng năng suất cá nhân, ghi chú, xử lý tài liệu nội bộ \\
\bottomrule
\end{tabularx}
\end{table}

\noindent\textbf{Kịch bản thực hành đồng hành}: Tệp \texttt{state\_machine\_contract\_demo.ts} (Chương 5) cài đặt Schema-First Zod, Máy trạng thái 4 bước và lớp \texttt{AICircuitBreaker}.

% CHƯƠNG 6
\chapter{Nguyên Lý Thiết Kế UX/UI Chống ``Nhựa'' (Anti-AI-Slop): Trải Nghiệm Chuẩn Nghệ Nhân}

\section{Bóc trần 5 căn bệnh của thẩm mỹ AI Slop}
Giao diện do AI tự sinh thường rất ``nhựa'' do mắc phải:
\begin{enumerate}[label={\arabic*.}]
    \item \textbf{Purple Gradient Craze}: Lạm dụng nền tím neon và dải màu cầu vồng vô nghĩa.
    \item \textbf{Cardocalypse}: Bọc mọi khối thông tin vào thẻ Card bo tròn quá đà và bóng đổ lơ lửng.
    \item \textbf{Iconitis}: Gắn các emoji và biểu tượng trang trí sáo rỗng vào trước mọi tiêu đề và nút bấm.
    \item \textbf{Hoạt họa co giật}: Áp dụng hiệu ứng chuyển động rung lắc làm rối mắt người dùng.
    \item \textbf{Văn phong phóng đại}: Dùng từ ngữ kịch tính hóa trong thông báo trạng thái.
\end{enumerate}

\section{Quy tắc phối màu 60--30--10 và Nhịp điệu lưới 8px}
\begin{itemize}[leftmargin=1.5em]
    \item \textbf{60\%}: Nền trung tính tĩnh lặng (Trắng hoặc Đen kẽm \texttt{\#09090B}).
    \item \textbf{30\%}: Bề mặt phẳng và đường phân cách mỏng 1px màu xám khói (\texttt{\#27272A}).
    \item \textbf{10\%}: Một màu nhấn chức năng duy nhất (Xanh cobalt hoặc ngọc). Cấm tuyệt đối dải màu gradient lòe loẹt.
    \item \textbf{Hệ thống lưới 8-Point Grid}: Mọi khoảng cách lề, đệm và chiều cao dòng đều là bội số của 8px (8, 16, 24, 32, 48, 64px) tạo nhịp thở thoáng đãng cho giao diện.
\end{itemize}

\section{Phân cấp kiểu chữ và Công thái học tương tác}
Sử dụng các họ font trung tính hiện đại (\textbf{Inter, Geist Sans, System UI}), tối đa 3 kích thước font và loại bỏ hoàn toàn chữ in hoa toàn bộ. Hỗ trợ phím tắt điều hướng nhanh (\texttt{Cmd+K}) và bắt buộc dùng \textbf{Skeleton Loader} mô phỏng chuẩn xác kích thước nội dung thay cho vòng xoay vô hồn.

\section{Cây quyết định: Lựa chọn phong cách thiết kế UI/UX}

\begin{table}[htbp]
\centering
\small
\caption{Cây quyết định định hình phong cách thiết kế giao diện}
\vspace{0.5em}
\begin{tabularx}{\textwidth}{L{3.2cm} L{3.8cm} L{4.5cm} Y}
\toprule
\textbf{Phong cách thiết kế} & \textbf{Đặc trưng nhận diện} & \textbf{Bộ màu tiêu chuẩn} & \textbf{Đối tượng \& Sản phẩm tối ưu} \\
\midrule
\textbf{Phong cách 1: Artisan Minimalist} & Phẳng, tối giản, viền 1px mỏng kẽm, font Inter/Geist Sans & Nền kẽm \texttt{\#09090B}, viền \texttt{\#27272A}, chữ trắng kẽm & Công cụ lập trình viên, sản phẩm SaaS hiện đại, bảng điều khiển \\
\textbf{Phong cách 2: Cyberpunk Dark Terminal} & Đen sâu, font monospace (JetBrains Mono), điểm nhấn xanh ngọc/amber & Nền đen tuyền \texttt{\#000000}, chữ xanh lục \texttt{\#10B981} hoặc cam & Công cụ dòng lệnh, bảo mật, DevOps, hacker toolkits \\
\textbf{Phong cách 3: Clean Enterprise SaaS} & Trắng sáng, nhã nhặn, tương phản chuẩn mực WCAG 2.2 AAA & Nền trắng tuyền \texttt{\#FFFFFF}, viền \texttt{\#E4E4E7}, nhấn xanh navy & Ứng dụng B2B cho doanh nghiệp, kế toán, quản trị nhân sự \\
\bottomrule
\end{tabularx}
\end{table}

\noindent\textbf{Minh họa thực hành}: Tệp \texttt{anti\_slop\_design\_system\_demo.html} (Chương 6) cung cấp trang so sánh trực quan đối đầu giữa giao diện AI Slop và giao diện Artisan UI.

% CHƯƠNG 7
\chapter{Chuyên Sâu Scripting, OS Internals \& Quy Trình Tự Chủ Vòng Lặp Kép}

\section{Bằng Chứng Thực Nghiệm 2026: Tại Sao Shell/CLI Đánh Bại Mọi Công Cụ Bọc Sẵn?}
Nghiên cứu mang tính bước ngoặt của nhóm tác giả \textit{Mak và cộng sự (arXiv:2609.11999, 2026)} mang tiêu đề: \textit{``Is Bash All You Need? An Empirical Study of Tool Interfaces for Enterprise Digital Worker Agents''} đã công bố một phát hiện thực nghiệm chấn động:
\begin{center}
\fbox{\parbox{0.88\textwidth}{
\textbf{Kết luận thực nghiệm từ arXiv:2609.11999 (2026):}\\
\textit{Các tác tử AI sử dụng giao diện dòng lệnh thuần túy (Shell/Bash) đạt tỷ lệ hoàn thành tác vụ cao hơn 38\% và tiết kiệm 42\% token ngữ cảnh so với việc sử dụng các công cụ bọc API chuyên biệt. Khả năng ghép nối dữ liệu qua đường ống (piping) cho phép AI giải quyết các tình huống biến đổi linh hoạt mà không bị kẹt trong các định dạng tĩnh.}
}}
\end{center}

\section{Linux Hacker Scripting: Nghệ Thuật Đường Ống Unix và IPC}
Trong triết lý Unix, mọi thứ đều là tệp (\textit{Everything is a file}), và sức mạnh tối thượng đến từ việc xâu chuỗi các chương trình nhỏ đơn nhiệm:
\begin{enumerate}[label={\arabic*.}]
    \item \textbf{File Descriptors (FD)}: Chuyển hướng cả stdout và stderr: \texttt{command > run.log 2>\&1} (hoặc \texttt{command \&> run.log}).
    \item \textbf{Named Pipe (FIFO)}: Giao tiếp liên tiến trình bất đồng bộ qua tệp ảo:
    \begin{lstlisting}[language=bash, caption={Co che Named Pipe FIFO va mo File Descriptor 3 tranh nghan EOF}]
# Khoi tao duong ong FIFO va mo FD 3
mkfifo /tmp/agent_ipc.pipe
exec 3<>/tmp/agent_ipc.pipe
# Doc su kien bat dong bo
while read -r event <&3; do
    echo "Xu ly su kien: $event"
done
    \end{lstlisting}
    \item \textbf{Signal Trapping}: Bắt các tín hiệu \texttt{SIGINT}, \texttt{SIGTERM} bằng lệnh \texttt{trap cleanup EXIT} để tự dọn dẹp tiến trình con và tệp khóa tạm.
    \item \textbf{Xử lý thời gian thực}: Dùng \texttt{inotifywait} làm Hot-Reload Daemon và \texttt{xargs -P} để chạy tác vụ song song trên đa nhân CPU.
\end{enumerate}

\section{Windows Hacker Scripting: Sức Mạnh PowerShell \& Object Pipeline}
Khác với Linux (dòng chữ thô), PowerShell trên Windows là một đường ống đối tượng (.NET Object Pipeline):
\begin{itemize}[leftmargin=1.5em]
    \item \textbf{Lọc đối tượng có cấu trúc}: Dùng \texttt{Where-Object \{ \$\_.WorkingSet -gt 200MB \}} truy cập trực tiếp RAM của tiến trình mà không cần regex.
    \item \textbf{P/Invoke Win32 API}: Nhúng trực tiếp C\# vào PowerShell để gọi các hàm hệ thống từ \texttt{user32.dll}.
    \item \textbf{AutoHotkey (AHK)}: Tự động hóa bàn phím và chuột cấp thấp, bắt Windows Message Hooks để điều khiển giao diện desktop tự động.
\end{itemize}

\section{Vòng Lặp Tự Phục Hồi Khép Kín (Self-Healing Loop-until-green)}
Bằng cách kích hoạt kịch bản \texttt{self\_healing\_test\_runner.py}, AI lắng nghe 4 kênh lỗi (Console, Hydration Exceptions, HTTP 4xx/5xx, Network failures), tự động xuất tệp \texttt{agent\_feedback.json} định vị chính xác URL và dòng mã nguồn gây lỗi kèm Exit Code 1. AI Agent đọc tệp này, tự sửa mã và chỉ kết thúc lượt chạy khi bài test trả về \textbf{Exit Code 0}.

\section{Cây quyết định: Lựa chọn môi trường tự động hóa \& Công nghệ kịch bản}

\begin{table}[htbp]
\centering
\small
\caption{Cây quyết định lựa chọn công nghệ tự động hóa \& kịch bản hệ thống}
\vspace{0.5em}
\begin{tabularx}{\textwidth}{L{3.2cm} L{3.8cm} L{4.5cm} Y}
\toprule
\textbf{Môi trường / Công nghệ} & \textbf{Cơ chế hoạt động} & \textbf{Vũ khí tối thượng} & \textbf{Bài toán áp dụng tối ưu} \\
\midrule
\textbf{Linux POSIX Bash} & Text streams, File Descriptors, Pipes, Signals & Named Pipes (FIFO), \texttt{inotifywait}, \texttt{xargs -P}, \texttt{trap} & Daemon ngầm, CI/CD, Container, tự động hóa build/test \\
\textbf{Windows PowerShell} & .NET Typed Object Pipeline, WMI/CIM, P/Invoke & \texttt{Get-Process}, \texttt{Add-Type} C\#, COM Automation & Quản trị hệ thống Windows, can thiệp Win32 Native API \\
\textbf{AutoHotkey (AHK)} & Windows Message Hook, gửi phím \& chuột cấp thấp & \texttt{WinActivate}, \texttt{ControlClick}, phím tắt toàn cục & Tự động hóa ứng dụng Desktop cũ không có API hoặc CLI \\
\textbf{Android ADB / scrcpy} & Android Debug Bridge, Framebuffer, UIAutomator XML & \texttt{adb exec-out screencap}, \texttt{uiautomator dump} & Thao tác ứng dụng di động Android, test bot mobile \\
\textbf{Python Hybrid Automation} & Thư viện chuẩn đa nền tảng, HTTPX, Subprocess, Pydantic & \texttt{subprocess.run}, \texttt{yt-dlp}, \texttt{ffmpeg-python} & Kịch bản đa nền tảng kết hợp xử lý media và gọi AI \\
\bottomrule
\end{tabularx}
\end{table}

\noindent\textbf{Bộ kịch bản thực hành đồng hành}: Thư mục Chương 7 trang bị sẵn \texttt{hacker\_linux\_daemon.sh}, \texttt{mobile\_adb\_controller.py}, \texttt{media\_transcriber\_pipeline.py}, và \texttt{office\_headless\_converter.sh}.

% PHỤ LỤC
\chapter*{Phụ Lục: Danh Mục Các Công Trình Nghiên Cứu 2026 Đã Thẩm Định Bằng Playwright}
\addcontentsline{toc}{chapter}{Phụ Lục: Danh Mục Nghiên Cứu Thực Nghiệm 2026}

Toàn bộ các luận điểm và số liệu trong giáo trình được xây dựng dựa trên 28 bài báo khoa học thực nghiệm mới nhất năm 2026 từ arXiv, đã được thẩm định tự động bằng kịch bản Playwright:

\begin{enumerate}[label={[\arabic*]}]
    \item \textbf{arXiv:2609.11999} -- \textit{Is Bash All You Need? An Empirical Study of Tool Interfaces for Enterprise Digital Worker Agents} (Khảo sát chứng minh Shell/CLI vượt trội 38\% so với API wrappers).
    \item \textbf{arXiv:2610.16460} -- \textit{The Interpreter as Developer: Building Local-First Tools through Expert-Led Vibe Coding} (Mô hình lắp ghép công cụ nguồn mở cục bộ).
    \item \textbf{arXiv:2511.19477} -- \textit{Building Browser Agents: Architecture, Security, and Practical Solutions} (Kiến trúc tác tử tự động hóa trình duyệt).
    \item \textbf{arXiv:2605.05509} -- \textit{WAAA! Web Adversaries Against Agentic Browsers} (Bảo mật cho browser agents).
    \item \textbf{arXiv:2610.14766} -- \textit{Last Secure Code Benchmark: Evaluating LLM Vulnerabilities in Functional Code Generation} (Khảo sát 65\% code AI dính lỗi bảo mật CWE).
    \item \textbf{arXiv:2609.09022} -- \textit{Cognitive Debt Accumulation in AI-Assisted Software Development} (Hiện tượng nợ nhận thức và suy thoái cơ sở mã).
    \item \textbf{arXiv:2609.29410} -- \textit{The Reimplementation Trap: Redundant Utility Generation in Agentic Codebases} (Bẫy tái triển khai mã trùng lặp của LLM).
\end{enumerate}

\end{document}
